\subsection{O-011: an exact multiplicity-one family detected by Hessians}

\paragraph{Finding O-011.}
Let $n>m\geq2$, and, for $1\leq d\leq n^2$, put
\[
              \lambda_d=(d(n-2)+2,2^{d-1})\vdash nd.
\]
Then the actual closure and ambient multiplicities are
\[
 K_n(\lambda_d)=a_{\lambda_d}(d[n])=1,\qquad
 q^m_{\lambda_d}(d[n])=
 \begin{cases}1&d\leq m^2,\\0&m^2<d\leq n^2.\end{cases}       \tag{O-011.1}
\]
In particular this entire family, including outputs with as many
as $m^2$ rows and positive permanent multiplicity, has no actual
multiplicity obstruction. No bound such as $n\geq m^{25}$ or
the sufficient weight-box inequality is needed for this family.

\begin{proof}
For a form $F$ of degree $n$ on $V$, define the degree-$d$
polynomial function
\[
 W_d(F)=\det\left(\frac{\partial^2F}{\partial x_i\partial x_j}
                                  (e_1)\right)_{1\leq i,j\leq d}.
\]
For an upper triangular $g\in\mathrm{GL}(V)$, write $g_d$ for
its leading $d\times d$ block. Since $ge_1=g_{11}e_1$, the chain
rule and homogeneity give
\[
 W_d(g^{-1}\cdot F)
       =g_{11}^{d(n-2)}\det(g_d)^2 W_d(F).
\]
Here $g^{-1}\cdot F=F\circ g$, consistently with the action on
the coordinate ring. Hence $W_d$ is a highest-weight vector of
weight $\lambda_d$, provided it is nonzero. The ambient
multiplicity is one by F-011 with $k=2$: in the Cauchy
decomposition on $d$ columns, the relevant multiplicity module is
\[
 S_{\lambda_d}\mathbb C^d
       =\det^2\otimes\operatorname{Sym}^{d(n-2)}\mathbb C^d.
\]
Its column weight $(n,\ldots,n)$ is the one-dimensional space
spanned, after removing the determinant factor, by
$\prod_i z_i^{n-2}$, and column permutations fix it. The case
$d=1$ is the same calculation directly.

At the identity matrix, the Hessian of $\det_n$ is the bilinear form
\[
               (A,B)\longmapsto\operatorname{tr}(A)\operatorname{tr}(B)
                                      -\operatorname{tr}(AB).
\]
It is nondegenerate: the trace-zero subspace is nondegenerate for
$-\operatorname{tr}(AB)$, its orthogonal complement is the span
of $I_n$, and the latter has square $n(n-1)\ne0$.
A nondegenerate symmetric bilinear form and a vector with nonzero
square admit a nondegenerate $d$-dimensional subspace containing
that vector for every $1\leq d\leq n^2$; diagonalize its
orthogonal complement to see this. Choose coordinates taking
$e_1$ to $I_n$ and the first $d$ coordinate directions to such
a subspace. Then $W_d$ is nonzero on a translate of $\det_n$.
Thus $K_n(\lambda_d)\geq1$, and the ambient multiplicity-one
bound gives $K_n(\lambda_d)=1$.

It remains to prove nonvanishing on the internally padded permanent.
Set $s=n-m>0$, $P=\operatorname{per}_m$, and $F=X_{11}^sP$.
Work first at the all-ones matrix. Put $h=(m-2)!$, let $u$ be
the all-ones vector of length $m^2$, let $e$ be the coordinate
vector for the $(1,1)$ entry, and put
\[
 B=(J_m-I_m)\otimes(J_m-I_m).
\]
Direct differentiation of the permanent gives
$P=m!$, $\nabla P=(m-1)!u$, and $H_P=hB$ at that point.
Therefore
\[
 \frac{H_F}{h}
   =B+s(m-1)(eu^T+ue^T)+s(s-1)m(m-1)ee^T.              \tag{O-011.2}
\]
The eigenvalues of $B$ are $(m-1)^2$, $-(m-1)$ and $1$, with
multiplicities $1$, $2(m-1)$ and $(m-1)^2$, respectively, so it
is invertible. Its inverse gives the following $2\times2$ Gram matrix:
\[
 G=\begin{pmatrix}e^TB^{-1}e&e^TB^{-1}u\\
                  u^TB^{-1}e&u^TB^{-1}u\end{pmatrix}
   =\frac1{(m-1)^2}
       \begin{pmatrix}(m-2)^2&1\\1&m^2\end{pmatrix}.
\]
Apply the matrix determinant lemma to (O-011.2), with
\[
 C=\begin{pmatrix}s(s-1)m(m-1)&s(m-1)\\s(m-1)&0\end{pmatrix}.
\]
Expanding $\det(I_2+CG)$ and factoring yields the exact identity
\[
 \frac{\det H_F}{\det H_P}
  =\frac{m+s-1}{m-1}\bigl(1-(m^2-3m+1)s\bigr)
  =\frac{n-1}{m-1}\bigl(1-(m^2-3m+1)(n-m)\bigr).       \tag{O-011.3}
\]
For positive integral $s$ this is nonzero except when $(m,n)=(3,4)$:
the integer $m^2-3m+1$ is $-1$ for $m=2$, $1$ for $m=3$,
and at least $5$ for $m\geq4$.

The exceptional point does not cause an exceptional coefficient.
For $m=3,n=4$, evaluate instead at the matrix with $X_{11}=2$
and every other entry equal to one. Then $F=16$. In the
orthonormal basis of sums and differences under interchange of the
last two rows and last two columns, the Hessian has blocks
\[
 A=\begin{pmatrix}
 4&2\sqrt2&2\sqrt2&10\\
 2\sqrt2&0&4&2\sqrt2\\
 2\sqrt2&4&0&2\sqrt2\\
 10&2\sqrt2&2\sqrt2&4
 \end{pmatrix},\qquad
 C_0=\begin{pmatrix}0&-2\sqrt2\\-2\sqrt2&-4\end{pmatrix}
 \quad\text{twice},\qquad (4).
\]
For clarity the four invariant vectors for the first block are
$e_{11}$, $(e_{12}+e_{13})/\sqrt2$,
$(e_{21}+e_{31})/\sqrt2$, and
$(e_{22}+e_{23}+e_{32}+e_{33})/2$; replacing sums by the
row or column differences gives the two sign blocks and the
double-sign line. Direct differentiation gives the displayed
blocks. Their determinants are $576,-8,-8,4$, so
$\det H_F=147456\ne0$.

In every case we have found a point $v$ in the $m^2$-dimensional
permanent variable space with $F(v)\ne0$ and nondegenerate Hessian.
Euler's identity gives
\[
                          H_F(v,v)=n(n-1)F(v)\ne0.
\]
As for the determinant, choose any $d\leq m^2$ and a
nondegenerate $d$-subspace containing $v$, and then an invertible
change of ambient coordinates. This gives a translate of $F$
on which $W_d$ is nonzero. Hence $q^m_{\lambda_d}(d[n])\geq1$;
the ambient bound gives equality. For $d>m^2$, the length
$\ell(\lambda_d)=d$ violates the necessary internal-padding
length condition (IP Theorem~1.8), so $q^m_{\lambda_d}(d[n])=0$.
This proves every case of (O-011.1).
\end{proof}

The construction is the determinant of a restriction of the second
jet of the universal degree-$n$ section at a chosen point. It is
equivariant under the Borel preserving that point and flag; its
weight is the displayed determinant-line character. Its derivatives,
matrix entries and identities before the convenient orthonormal
change of basis are all defined over $\mathbb Z$. The multiplicity
interpretation uses the characteristic-zero Schur decomposition.
This places the explicit calculation in the equivariant locally
free-sheaf and Schur-functor dictionary of SGA~3, I, 4.7.2--4.7.3,
SGA~6, VI, Theorem~3.3, and cote~133, scan~028, without asserting
the same semisimple multiplicities in residue characteristics.
